\section{The LASR-MCP Architecture}
\label{sec:architecture}

\begin{figure*}[t]
\centering
\resizebox{0.90\textwidth}{!}{%
\begin{tikzpicture}[
    node distance=1.4cm and 1.2cm,
    block/.style={rectangle, draw=black, thick, fill=blue!5, text width=2.3cm, text centered, rounded corners=2pt, minimum height=1.1cm, font=\footnotesize},
    process/.style={rectangle, draw=black, thick, fill=orange!10, text width=2.5cm, text centered, rounded corners=2pt, minimum height=1.1cm, font=\footnotesize},
    decision/.style={diamond, draw=black, thick, fill=green!10, text width=1.4cm, text centered, aspect=1.8, font=\footnotesize},
    output/.style={rectangle, draw=black, thick, fill=teal!10, text width=2.3cm, text centered, rounded corners=2pt, minimum height=1.1cm, font=\footnotesize},
    line/.style={draw=black, -stealth, thick},
    dashline/.style={draw=black, -stealth, thick, dashed}
]

% Nodes
\node [block] (input) {\textbf{Tool Contract Tree}\\ $\mathcal{T} = (\mathcal{V}, \mathcal{E}, \tau)$ \\ \scriptsize Exhaustive Schema (552 tok)};

\node [process, right=of input] (bundler) {\textbf{Priority Bundler}\\ Lean Projection $C_I^* \cup \hat{C}_T(x)$\\ Upfront $k=2$ ($S_0=359.4$ tok)};

\node [block, right=of bundler] (llm) {\textbf{LLM Execution}\\ \scriptsize\texttt{gemini-3.5-}\\ \scriptsize\texttt{flash-lite}\\ Attempt $t$};

\node [decision, right=of llm] (val) {\textbf{Contract}\\ \textbf{Valid?}};

\node [output, right=of val] (success) {\textbf{Execution Pass}\\ Validated Payload\\ $y_t \models S$};

\node [process, below=1.2cm of llm] (slicer) {\textbf{Reactive Slicer}\\ Multi-Error Delta $\Delta S$\\ $\mathcal{T}_{t+1} = \mathcal{T}_t \cup \text{Slice}(\mathcal{T}, \mathcal{E}_t)$\\ \scriptsize Mean $\Delta S = 45.2$ tok};

% Paths
\path [line] (input) -- (bundler);
\path [line] (bundler) -- (llm);
\path [line] (llm) -- (val);
\path [line] (val) -- node[above, font=\footnotesize\bfseries] {Yes} (success);
\path [line] (val) |- node[right, font=\footnotesize\bfseries, near start] {No (Errors $\mathcal{E}_t$)} (slicer);
\path [line] (slicer) -- node[left, font=\footnotesize] {Inject $\Delta S$} (llm);

% Background box for in-tree recovery
\begin{scope}[on background layer]
    \node [draw=blue!40, dashed, fill=blue!2, rounded corners=4pt, fit=(llm) (val) (slicer), label=above:{\footnotesize\textbf{Dynamic In-Tree Recovery Loop} \scriptsize(Bound: $K \le M - |A_0|)$}] {};
\end{scope}

\end{tikzpicture}%
}
\caption{System Architecture of LASR-MCP: The full tool contract tree $\mathcal{T}$ is structurally projected into an invariant foundation $C_I^*$ with $k=2$ prioritized constraint bundles ($S_0 = 359.4$ tokens). Upon validation failure, the reactive slicer dynamically splices concurrent violated constraint subtrees ($\Delta S = 45.2$ tokens) corresponding to the diagnostic error set $\mathcal{E}_t$ into Attempt $t+1$, bounded by the monotonic convergence invariant $K \le M - |A_0|$.}
\label{fig:system_architecture}
\end{figure*}


To resolve the trade-off between bloated context overhead and brittle schema truncation, we introduce \textbf{Learned Adaptive Schema Resolution (LASR-MCP)}. Rather than treating tool schemas as static monolithic declarations or appending unanchored constraint slices, LASR-MCP formulates the active tool schema at execution attempt $t$ as a three-tier union:
\begin{equation}
\label{eq:three_tier}
C_{\text{active}}^{(t)} = C_I^* \;\cup\; \hat{C}_T(x) \;\cup\; C_R^{(t)}
\end{equation}
where $C_I^*$ represents the schema-wide \textit{Invariant Foundation}, $\hat{C}_T(x)$ is the \textit{Proactive Task-Conditioned Constraint Bundle} predicted from user intent $x$, and $C_R^{(t)}$ is the accumulated set of \textit{Dynamically Recovered In-Tree Constraints} derived from validation feedback up to round $t$.

\subsection{Tier 1: Invariant Foundation Floor ($C_I^*$)}
Empirical inspection of complex JSON Schema specifications (e.g., enterprise purchase order dispatching) reveals that tool execution cannot succeed on intent-specific properties alone. Certain structural constraints are \textit{invariants} required for all syntactically well-formed submissions regardless of user goal (e.g., supplier identification regex format, item array container specifications, and SKU patterns).

Let $\mathcal{C}_{\text{full}}$ denote the complete schema constraint set ($552$ tokens). We define the Invariant Foundation $C_I^*$ as the minimal constraint subset among tested candidates satisfying the pre-specified baseline structural-validity floor without user-intent conditioning:
\begin{equation}
C_I^* = \arg\min_{C \in \mathcal{C}_{\text{cand}}} |C| \quad \text{s.t.} \quad \mathcal{V}_{\text{struct}}(C) \ge \tau_{\text{floor}}
\end{equation}
where $\mathcal{C}_{\text{cand}}$ is the set of evaluated candidate invariant subsets and $\tau_{\text{floor}} = 0.70$ specifies the baseline validity floor, chosen \textit{a priori} as the minimum operational viability threshold to avoid retry cascade failure. As demonstrated in Section~\ref{sec:results}, $C_I^* = 339$ tokens is the empirically selected invariant foundation meeting this threshold among tested tiers for the evaluated purchase-order schema; compressing below this boundary triggers structural collapse in nested item definitions.

\subsection{Tier 2: Proactive Intent-Conditioned Predictor ($\hat{C}_T(x)$)}
Given a user query $x \in \mathcal{X}$, the agent's goal is to predict the minimal subset of task-conditioned constraints $\hat{C}_T(x) \subseteq \mathcal{C}_{\text{task}}$ necessary for successful execution. We formulate this as a multi-label classification problem over the task-conditioned label space $\mathcal{L} = \{c_1, c_2, \dots, c_m\}$ (e.g., specific warehouse routing enums, order priority restrictions, and budget bounds).

We employ a frozen linear semantic predictor:
\begin{equation}
\hat{y}_i = \sigma\left( \mathbf{w}_i^T \phi(x) + b_i \right), \quad i \in \{1, \dots, m\}
\end{equation}
where $\phi(x)$ denotes normalized term-salience features extracted from the query. The decision threshold $\theta$ was calibrated strictly on the 48 training tasks and frozen prior to evaluation. To construct the proactive bundle of size up to $k$, the highest-ranked constraints exceeding $\theta$ (up to $k$) are selected:
\begin{equation}
\hat{C}_T(x) = \text{Top-}k \left( \{ c_i \in \mathcal{L} \mid \hat{y}_i \ge \theta \} \right)
\end{equation}
We formally denote by $m_{\text{task}}(x) = |C_T^*(x)|$ the number of required non-invariant task-conditioned constraints demanded by user intent $x$. The condition $m_{\text{task}}(x) \le k$ is an operational regime descriptor characterizing when task requirements do not exceed the upfront budget capacity $k$, rather than a guarantee that the heuristic predictor selects all required constraints on Attempt 1 (e.g., in Cluster $C_1$, first-pass hierarchical success is $75.0\%$, though final structural success recovers to $100.0\%$). Conversely, $m_{\text{task}}(x) > k$ defines regimes where omissions are structurally guaranteed.

\subsection{Tier 3: Dynamic In-Tree Recomposition ($C_R^{(t)}$)}
\label{sec:recomposition}
When an unpredicted constraint is omitted from $C_I^* \cup \hat{C}_T(x)$, the LLM may submit arguments $y_t$ that fail JSON Schema validation. A simple agent implementation may return raw validator error strings into conversation history or retry without modifying the exposed schema. 

LASR-MCP introduces \textit{in-tree reactive recomposition}. Reactive recovery operates as pass-level multi-error constraint expansion rather than single-error sequential repair. Our validator is configured to return all observed violations $\text{Errors}(y_t) = \{e_1, e_2, \dots\}$ in a single validation pass, rather than short-circuiting on the first failure. Analogous to an enterprise validator reporting all diagnostic issues across a payload, the reactive operator extracts all concurrent violations simultaneously:
\begin{equation}
\label{eq:reactive_op}
C_R^{(t+1)} = C_R^{(t)} \cup \bigcup_{e \in \text{Errors}(y_t)} \left( \text{Fragment}(e) \cap \mathcal{S}_{\text{specs}} \setminus C_{\text{active}}^{(t)} \right)
\end{equation}
where $\text{Fragment}(e)$ deterministically maps the RFC 6901 JSON Pointer error path~\cite{rfc6901} and validator-reported JSON Schema keyword (e.g., \texttt{path='/warehouse\_code'}, \texttt{keyword='enum'}) to the canonical constraint node $\Delta_e \in \mathcal{S}_{\text{specs}}$.

The updated schema specification $S_{t+1}$ is then structurally spliced into the native JSON Schema hierarchy:
\begin{equation}
S_{t+1} = \mathcal{R}\left( S_I, C_{\text{active}}^{(t+1)} \right)
\end{equation}
where $\mathcal{R}$ is the in-tree recomposition function that inserts properties, pattern regexes, and enums directly into their respective object trees without creating duplicate definition nodes or schema syntax errors.

\subsection{Incremental Token Accounting Metrics}
To evaluate token economics, we audit token consumption at every turn:
\begin{itemize}
    \item \textbf{Initial Schema Overhead ($S_0$)}: Tokens required to serialize $S_{\text{active}}^{(1)}$.
    \item \textbf{Incremental Expansion ($\Delta S_t$)}: Additional schema tokens introduced at attempt $t \ge 2$: $\Delta S_t = \text{Tok}(S_t) - \text{Tok}(S_{t-1})$.
    \item \textbf{Structural Tokens-to-Success ($\text{TTS}_{\text{struct}}$)}: Cumulative recorded prompt-plus-candidate tokens across all attempts conditional on structurally successful trajectories ($y_{T_{\text{succ}}} \models \mathcal{S}$ within attempt budget $T_{\text{succ}} \le 3$):
    \begin{equation}
    \begin{split}
    TTS_{\text{struct}} &= \sum_{t=1}^{T_{\text{succ}}} \text{Recorded } (P+C)_t \\
    &= \sum_{t=1}^{T_{\text{succ}}} \bigl( \texttt{promptTokenCount}_t \\
    &\qquad\quad + \texttt{candidatesTokenCount}_t \bigr)
    \end{split}
    \end{equation}
    where provider-recorded prompt tokens natively encompass both non-schema text ($P_t$) and the serialized active schema ($\text{Tok}(S_t)$), preventing offline double-counting, while $C_t$ represents candidate completion tokens.
    \item \textbf{Total Workload Tokens}: Sum of all tokens consumed across all attempts up to the three-attempt budget, including failed retry loops. We explicitly distinguish total workload on the historical development benchmark (32 unique task trials per arm) from the canonical live benchmark (640 task executions per arm across 20 fixed-seed repetitions; 3,840 total trials across all six arms).
    \item \textbf{Structural Recovery Rate}: Fraction of initially structurally invalid trajectories that subsequently satisfy JSON Schema validation within the attempt budget.
\end{itemize}
For workload-token reporting, we use the provider-recorded prompt-plus-candidate token subtotal per request ($\text{Recorded } P+C = \texttt{promptTokenCount} + \texttt{candidatesTokenCount}$, which identically equals provider \texttt{totalTokenCount} across all 7,860 response-bearing calls in our benchmark dataset as thinking and client-side tool prompt tokens are zero; we report this as an audited property of our experimental telemetry rather than an invariant across all API features). Serialized schema tokens ($S_0, \Delta S$) are audited separately as structural diagnostics without double-counting them into provider prompt totals.

\subsection{Formal Foundation \& Invariant Termination Bound}
\label{sec:formal_proof}

Let $\mathcal{S}$ denote a canonical tool specification expressed in JSON Schema Draft 2020-12, represented as an abstract syntax tree (AST). We define $\mathcal{C} = \{c_1, c_2, \dots, c_M\}$ as the finite universe of structural constraint nodes in $\mathcal{S}$ ($|\mathcal{C}| = M = M_{\text{schema}} < \infty$), where each constraint node $c_i$ corresponds to an unambiguous validation keyword (e.g., \texttt{pattern}, \texttt{enum}, \texttt{minimum}, \texttt{required}) situated at a specific JSON Pointer path. We explicitly distinguish this total schema universe size $M_{\text{schema}}$ from task-conditioned constraint requirements $m_{\text{task}}(x)$.

The reactive framework operates over a sequence of active constraint sets $A_t \subseteq \mathcal{C}$ at recovery step $t \ge 0$, initialized as $A_0 = C_I^* \cup \widehat{C}_T(x)$. At step $t$, the generator outputs argument payload $y_t$. If $y_t$ violates schema constraints, the validator returns diagnostics $\mathcal{E}(y_t, \mathcal{S}) \neq \emptyset$. The deterministic slicer $\sigma$ maps each diagnostic to its canonical schema constraint node. To ensure strict non-redundant growth, the newly identified constraint delta is defined by explicitly filtering previously active nodes:
\begin{equation}
\Delta_t = \left( \bigcup_{e \in \mathcal{E}(y_t, \mathcal{S})} \sigma(e, \mathcal{S}) \right) \setminus A_t
\end{equation}
The reactive recomposer then updates the active constraint set in-tree: $A_{t+1} = A_t \cup \Delta_t$.

\begin{definition}[Admissible Reactive System]
A reactive tool-calling system $(\mathcal{S}, \sigma, \mathcal{G})$ is admissible if:
\begin{enumerate}
    \item \textbf{Finite Constraint Universe}: The schema AST contains a finite number of constraint nodes, $|\mathcal{C}| = M = M_{\text{schema}} < \infty$.
    \item \textbf{Non-Empty Fault Localization}: On any validation failure at step $t$, the slicer identifies at least one previously inactive constraint: $|\Delta_t| \ge 1$ and $\Delta_t \cap A_t = \emptyset$.
    \item \textbf{Strict Monotonicity}: Each localized constraint is merged idempotently without pruning active constraints, ensuring $A_t \subset A_{t+1} \subseteq \mathcal{C}$.
\end{enumerate}
\end{definition}

\begin{theorem}[Finite Monotonic Recovery Bound]
\label{thm:termination}
Let $(\mathcal{S}, \sigma, \mathcal{G})$ be an admissible reactive tool-calling system with initial active constraint set $A_0 \subseteq \mathcal{C}$ and full schema universe $M_{\text{schema}} = |\mathcal{C}| = M$. The dynamic in-tree recovery process terminates in at most:
\begin{equation}
\boxed{K \le M - |A_0|}
\end{equation}
schema-expansion transitions (where $M \equiv M_{\text{schema}}$), resulting in either a structurally valid payload or complete exhaustion of the schema ($A_K = \mathcal{C}$).
\end{theorem}

\begin{IEEEproof}
By the non-empty fault localization condition, whenever $y_t$ is invalid, $\Delta_t \subseteq \mathcal{C} \setminus A_t$ and $|\Delta_t| \ge 1$. Consequently:
\begin{equation}
|A_{t+1}| = |A_t \cup \Delta_t| = |A_t| + |\Delta_t| \ge |A_t| + 1
\end{equation}
Thus, the active cardinality sequence is strictly monotonically increasing: $|A_0| < |A_1| < \dots < |A_K| \le M$. By induction across $K$ transitions, $|A_K| \ge |A_0| + K$. Because $|A_K| \le M$, it follows directly that $|A_0| + K \le M \implies K \le M - |A_0|$. Because $\mathcal{C}$ is finite and omitted constraints strictly decrease, the recovery loop cannot cycle or run indefinitely, guaranteeing termination in at most $M - |A_0|$ transitions.
\end{IEEEproof}
We emphasize that Theorem~\ref{thm:termination} bounds dynamic schema expansion under its stated admissibility conditions; it guarantees that the active schema contract monotonically expands to full specification in at most $M - |A_0|$ transitions, but does not guarantee that a stochastic model will always generate a valid payload if the model violates already-exposed constraints or if execution is constrained by an operational attempt budget (such as $T_{\text{max}} = 3$).

\begin{remark}[Analysis of the Invalid Tree-Depth Bound $K \le D$]
Earlier heuristic drafts proposed a bound of $K \le \min(D, M - |C_I^*|)$, where $D$ represents maximum AST depth. This claim does not hold in general. Consider two sibling leaf constraints at depth $D = 1$: $c_1 =$ \texttt{properties/warehouse\_code/enum} and $c_2 =$ \texttt{properties/order\_priority/enum}. On Attempt 1 ($A_0 = \emptyset$), the generator may violate $c_1$ (splicing $c_1$). On Attempt 2, it may satisfy $c_1$ but violate $c_2$ (splicing $c_2$), converging on Attempt 3. This yields $K = 2$ transitions even though $D = 1$. Since $2 \not\le 1$, tree depth does not constrain recovery passes. Theorem~\ref{thm:termination} establishes the correct cardinality-based bound.
\end{remark}

\begin{proposition}[Structural Validity Guarantee]
If the generator $\mathcal{G}$ satisfies the \emph{exposed-constraint compliance assumption}---namely, that $\mathcal{G}$ satisfies all constraints in $A_t$ and failures arise strictly from omitted constraints in $\mathcal{C} \setminus A_t$---then the recovery loop terminates with a structurally valid invocation ($y_K \models \mathcal{S}$) in $K \le M - |A_0|$ transitions.
\end{proposition}

\begin{remark}[Deterministic Schema Convergence vs.\ Stochastic Decoding Adherence]
It is critical to distinguish between \emph{deterministic schema state convergence} and \emph{stochastic LLM decoding adherence}. Theorem~\ref{thm:termination} guarantees that the schema contract itself strictly and monotonically expands to include all violated constraints within at most $M - |A_0|$ steps. However, once a constraint is exposed in the prompt, actual generative compliance depends on the model's stochastic decoding adherence. In our live $N=3,840$ empirical evaluation on \texttt{gemini-3.5-flash-lite}, this distinction accounts for the observed $96.9\%$ vs.\ $100.0\%$ final structural completion rate: the schema state monotonically reached full specification, but the model occasionally generated invalid edge-case syntax on Attempt 3 in complex multi-item nested arrays ($C_3$).
\end{remark}


